
\chapter{Fundamentos de los cálculos de estructura electrónica}
\label{ap:fundamentos-gaussian}
\section{Programas y conceptos generales}
Los estudios de estructura electrónica permiten relacionar la geometría de una molécula con su energía y con la distribución de sus electrones. En este trabajo, estos conceptos sirven de base para interpretar la caracterización del fulereno y de sus derivados con glicina. Gaussian 16 es el programa de cálculo considerado en el proyecto; GaussView es la interfaz utilizada para construir y visualizar las estructuras. El programa, el método teórico y el conjunto de bases cumplen funciones distintas y, por ello, deben identificarse por separado [\ref{anexo-ref-1}].

\section{Teoría del funcional de la densidad}
La teoría del funcional de la densidad, conocida como DFT, describe el estado electrónico fundamental a partir de la densidad electrónica $\rho(\mathbf r)$. Esta magnitud representa cómo se distribuyen los electrones en el espacio y su integral corresponde al número de electrones del sistema. En el esquema de Kohn--Sham, la energía se escribe como
\begin{equation}
 E[\rho]=T_s[\rho]+\int v_{\mathrm{ext}}(\mathbf r)\rho(\mathbf r)\,d\mathbf r
 +J[\rho]+E_{\mathrm{xc}}[\rho]+E_{\mathrm{NN}}.
\end{equation}
Aquí, $T_s$ es la energía cinética del sistema auxiliar de electrones no interactuantes; $v_{\mathrm{ext}}$ representa el potencial externo, que incluye la atracción de los núcleos; $J$ es la repulsión electrostática clásica entre electrones; $E_{\mathrm{xc}}$ reúne las contribuciones de intercambio y correlación, incluida la diferencia de energía cinética respecto al sistema auxiliar; y $E_{\mathrm{NN}}$ es la repulsión entre núcleos. La aproximación empleada para $E_{\mathrm{xc}}$ influye en las propiedades calculadas [\ref{anexo-ref-2}].

\section{Funcional B3LYP y conjunto de bases}
El nivel de teoría declarado en Materiales y métodos es B3LYP/6-311G(d,p). B3LYP es un funcional híbrido que combina una fracción de intercambio de Hartree--Fock con términos de intercambio y correlación dependientes de la densidad. Su denominación incorpora la formulación de tres parámetros de Becke y la correlación de Lee, Yang y Parr [\ref{anexo-ref-3}, \ref{anexo-ref-4}]. Se trata de una aproximación: su uso no garantiza, por sí solo, que todas las propiedades de una molécula se describan con la misma precisión.

El conjunto de bases proporciona las funciones con las que se representan los orbitales. En 6-311G(d,p), la valencia se describe mediante tres componentes de distinta extensión radial. Las funciones de polarización $d$, en los átomos distintos del hidrógeno, y $p$, en el hidrógeno, permiten representar cambios en la forma de la distribución electrónica [\ref{anexo-ref-5}]. Esta notación no incluye funciones difusas, que se indican mediante signos $+$. Por tanto, no debe confundirse con 6-311+G(d,p) ni con 6-311++G(d,p).

\section{Optimización, frecuencias y comparación energética}
La optimización geométrica busca una estructura estacionaria sobre la superficie de energía potencial. El análisis vibracional permite caracterizar ese punto mediante la curvatura de la superficie: un mínimo local presenta curvaturas positivas en sus modos vibracionales. Por ello, el protocolo declara la comprobación de ausencia de frecuencias imaginarias. Esta comprobación debe respaldarse con las salidas correspondientes y no demuestra que se haya localizado el mínimo global entre todas las conformaciones posibles.

Las energías relativas se expresan respecto a una referencia común:
\begin{equation}
 \Delta E_i=E_i-E_{\mathrm{ref}}.
\end{equation}
Si la referencia no es la estructura de menor energía, pueden aparecer valores negativos. Además, la comparación directa exige que los sistemas tengan la misma composición y que las energías se hayan obtenido con criterios equivalentes. La energía electrónica, la energía corregida por el punto cero, la entalpía y la energía libre de Gibbs son cantidades distintas; no deben presentarse bajo un mismo encabezado sin aclarar cuál se utilizó. Una diferencia entre isómeros tampoco constituye, por sí sola, una energía de formación.

\chapter{Funciones de Fukui y descriptores de reactividad}
\label{ap:reactividad}

La reactividad electrónica puede analizarse desde dos perspectivas. Los descriptores globales resumen la respuesta de toda la molécula ante cambios en su número de electrones. Los descriptores locales muestran cómo se distribuye esa respuesta en distintas regiones. Ambos enfoques se complementan, pero no sustituyen el estudio de una reacción ni demuestran actividad biológica.

\section{Funciones de Fukui en el espacio}
La función de Fukui describe el cambio de la densidad electrónica cuando varía el número de electrones, manteniendo fijo el potencial externo:
\begin{equation}
 f(\mathbf r)=\left(\frac{\partial\rho(\mathbf r)}{\partial N}\right)_{v_{\mathrm{ext}}}.
\end{equation}
En una molécula, esto implica conservar la disposición nuclear al comparar las densidades. La aproximación por diferencias finitas distingue tres respuestas [\ref{anexo-ref-6}]:
\begin{align}
 f^+(\mathbf r)&=\rho_{N+1}(\mathbf r)-\rho_N(\mathbf r),\\
 f^-(\mathbf r)&=\rho_N(\mathbf r)-\rho_{N-1}(\mathbf r),\\
 f^0(\mathbf r)&=\tfrac12\left[\rho_{N+1}(\mathbf r)-\rho_{N-1}(\mathbf r)\right].
\end{align}
La función $f^+$ se relaciona con la respuesta al recibir electrones y con regiones susceptibles a un ataque nucleofílico. La función $f^-$ describe la respuesta al perder electrones y se asocia con regiones susceptibles a un ataque electrófilo. La función $f^0$ es el promedio de ambas respuestas y se emplea en el análisis de reactividad radicalaria. El signo del superíndice identifica la variación electrónica considerada; no representa la carga formal de un átomo.

\section{Ionización y afinidad electrónica}
El potencial de ionización $I$ expresa el costo energético de retirar un electrón; la afinidad electrónica $A$ expresa la energía liberada al incorporarlo, con la convención de signo aquí adoptada. A partir de energías totales,
\begin{align}
 I&=E(N-1)-E(N),\\
 A&=E(N)-E(N+1).
\end{align}
Si las energías se comparan en una misma geometría, las magnitudes son verticales. Si cada estado se encuentra en su propia geometría de equilibrio, corresponden a diferencias adiabáticas, cuya definición debe precisar también las correcciones energéticas incluidas.

Las expresiones $I\approx-\varepsilon_{\mathrm{HOMO}}$ y $A\approx-\varepsilon_{\mathrm{LUMO}}$ constituyen aproximaciones orbitales. No son equivalentes a las diferencias de energía total. En DFT, la discontinuidad de la derivada de la energía limita la identificación directa de la afinidad electrónica con el negativo de la energía del LUMO del sistema neutro [\ref{anexo-ref-8}]. Por ello, los dos enfoques deben distinguirse al describir los resultados.

\section{Potencial químico, electronegatividad y dureza}
El potencial químico electrónico $\mu$ y la electronegatividad $\chi$ se aproximan mediante
\begin{equation}
 \mu\approx-\frac{I+A}{2},\qquad \chi=-\mu.
\end{equation}
Para evitar ambigüedades, en este apéndice se adopta la convención de dureza como curvatura, $\eta\approx I-A$. Con ella, la blandura global y el índice de electrofilicidad se expresan como
\begin{equation}
 S=\frac{1}{\eta},\qquad \omega=\frac{\mu^2}{2\eta}.
\end{equation}
La dureza describe la resistencia de la distribución electrónica a cambiar; la blandura expresa la respuesta inversa. La electrofilicidad combina el potencial químico y la dureza para describir la estabilización asociada con aceptar densidad electrónica en un modelo de respuesta cuadrática [\ref{anexo-ref-9}].

Existe también la convención de semidureza $\eta_{1/2}=(I-A)/2$. Bajo esa definición, las mismas cantidades anteriores serían $S=1/(2\eta_{1/2})$ y $\omega=\mu^2/(4\eta_{1/2})$. Por tanto, la fórmula y la convención deben declararse juntas antes de comparar valores de distintas fuentes. Si $I$ y $A$ se expresan en eV, $\mu$, $\chi$, $\eta$ y $\omega$ se expresan en eV, mientras que $S$ se expresa en eV$^{-1}$.

Estos descriptores se presentan como fundamento teórico. No se incorporan como resultados nuevos, pues el protocolo contiene valores de ionización y afinidad electrónica, pero no una tabla documentada de todos los índices anteriores.

\section{Potencial electrostático molecular}
El potencial electrostático molecular representa la interacción de una carga positiva de prueba con la distribución de núcleos y electrones. En unidades atómicas,
\begin{equation}
 V(\mathbf r)=\sum_A\frac{Z_A}{|\mathbf r-\mathbf R_A|}
 -\int\frac{\rho(\mathbf r')}{|\mathbf r-\mathbf r'|}\,d\mathbf r'.
\end{equation}
Su representación sobre una superficie molecular permite reconocer regiones de potencial positivo y negativo. Para comparar mapas es necesario conocer las unidades, la superficie representada y la escala de colores. El potencial electrostático no equivale a la densidad electrónica ni al momento dipolar. Asimismo, un potencial cercano a cero no demuestra que una región sea hidrofóbica.





\clearpage
\chapter*{Referencias de los apéndices}
\phantomsection
\addcontentsline{toc}{chapter}{Referencias de los apéndices}
Las referencias de los apéndices se numeran de manera independiente de la bibliografía del protocolo, según su primera aparición en los apéndices.
\begingroup
\renewcommand{\labelenumi}{[\arabic{enumi}]}
\begin{enumerate}
\setlength{\itemsep}{0.65\baselineskip}
\item\label{anexo-ref-1} M. J. Frisch et al. \textit{Gaussian 16}. Gaussian, Inc., 2016. La revisión específica del programa utilizado queda pendiente de corroborar.

\item\label{anexo-ref-2} W. Kohn y L. J. Sham. ``Self-Consistent Equations Including Exchange and Correlation Effects''. \textit{Physical Review}, 140, A1133--A1138, 1965. DOI: \href{https://doi.org/10.1103/PhysRev.140.A1133}{10.1103/PhysRev.140.A1133}.

\item\label{anexo-ref-3} A. D. Becke. ``Density-functional thermochemistry. III. The role of exact exchange''. \textit{The Journal of Chemical Physics}, 98, 5648--5652, 1993. DOI: \href{https://doi.org/10.1063/1.464913}{10.1063/1.464913}.

\item\label{anexo-ref-4} C. Lee, W. Yang y R. G. Parr. ``Development of the Colle-Salvetti correlation-energy formula into a functional of the electron density''. \textit{Physical Review B}, 37, 785--789, 1988. DOI: \href{https://doi.org/10.1103/PhysRevB.37.785}{10.1103/PhysRevB.37.785}.

\item\label{anexo-ref-5} R. Krishnan, J. S. Binkley, R. Seeger y J. A. Pople. ``Self-consistent molecular orbital methods. XX. A basis set for correlated wave functions''. \textit{The Journal of Chemical Physics}, 72, 650--654, 1980. DOI: \href{https://doi.org/10.1063/1.438955}{10.1063/1.438955}.

\item\label{anexo-ref-6} R. G. Parr y W. Yang. ``Density functional approach to the frontier-electron theory of chemical reactivity''. \textit{Journal of the American Chemical Society}, 106(14), 4049--4050, 1984. DOI: \href{https://doi.org/10.1021/ja00326a036}{10.1021/ja00326a036}.

\item\label{anexo-ref-7} W. Yang y W. J. Mortier. ``The use of global and local molecular parameters for the analysis of the gas-phase basicity of amines''. \textit{Journal of the American Chemical Society}, 108(19), 5708--5711, 1986. DOI: \href{https://doi.org/10.1021/ja00279a008}{10.1021/ja00279a008}.

\item\label{anexo-ref-8} J. P. Perdew, R. G. Parr, M. Levy y J. L. Balduz. ``Density-Functional Theory for Fractional Particle Number: Derivative Discontinuities of the Energy''. \textit{Physical Review Letters}, 49, 1691--1694, 1982. DOI: \href{https://doi.org/10.1103/PhysRevLett.49.1691}{10.1103/PhysRevLett.49.1691}.

\item\label{anexo-ref-9} R. G. Parr, L. von Szentpály y S. Liu. ``Electrophilicity Index''. \textit{Journal of the American Chemical Society}, 121(9), 1922--1924, 1999. DOI: \href{https://doi.org/10.1021/ja983494x}{10.1021/ja983494x}.

\end{enumerate}
\endgroup